\section{What the forecast is for}
\label{sec:sizing}

A predicted standard deviation is not a result on its own. Nobody trains a model in order to learn
how much its score would move under a different seed; they want to know whether one recipe beats
another, and the deviation is the thing that says how many runs settle it. This section measures
the forecast by that standard.

\subsection{The decision}

Two configurations are to be ordered. A practitioner trains each for a seed, looks at the lead,
and either calls it or trains another seed of each. Stopping when the lead $\bar D_k$ after $k$
seeds per arm clears $z$ standard errors,
\begin{equation}
  \label{eq:stop}
  |\bar D_k| \;\ge\; z\,\frac{\sigma_D}{\sqrt{k}},
  \qquad
  \sigma_D=\sqrt{\sigma_A^2+\sigma_B^2},
\end{equation}
is the rule this is an instance of. It needs $\sigma_D$, and that is the difficulty: at $k=1$
there is no empirical deviation to put there, because one seed per arm carries no information
about spread. The choices that remain are to fix $k$ in advance, which is what is normally done, or
to spend seeds building a variance estimate before the comparison can begin. Section~\ref{sec:seeds}
supplies $\sigma_A$ and $\sigma_B$ from the first run of each arm, so \eqref{eq:stop} is available
from $k=1$.

Note what \eqref{eq:stop} does \emph{not} require. Multiplying every $\sigma$ by a common factor is
undone by dividing $z$ by the same factor, so a forecast that is uniformly off by a constant
reaches the same operating points as one that is not, provided $z$ is chosen on a range wide enough
to allow the move. What matters is whether the forecast gets the \emph{relative} deviation right,
across configurations. We check both claims below.

\subsection{Setup}

The scores are the $\numSizeSeeds$ per-seed scores already stored for each of the
$\numSizeArms$ configurations of Section~\ref{sec:exp-forecast}; no new training is done, and every
policy sees the same seeds in the same order, so the comparison is paired. We compare
configurations that share a task, and keep the $\numSizePairs$ pairs whose ordering over all
$\numSizeSeeds$ seeds is unambiguous at four standard errors, so that a disagreement is a mistake
rather than a coin landing the other way up. Each pair is replayed under $400$ random seed
orderings, and a policy is a curve rather than a point, since $z$ trades runs against mistakes.
Reading every curve at a common mistake rate is what makes two policies comparable; we report the
runs each needs to hold its mistake rate at $\numSizeTarget\%$, and both arms are counted, so
$k$ seeds per arm is $2k$ runs.

Six policies are compared. \emph{Fixed budget} spends a predetermined $k$ and never stops early.
\emph{Deviation from the seeds} estimates $\sigma_D$ from the seeds in hand, and cannot stop before
$k=2$. \emph{Constant deviation} uses one number for every configuration, the median measured
deviation, which is the sharpest form of the assumption that variance does not depend on the
recipe. \emph{Frozen forecast} is the predictor of Section~\ref{sec:estimation} available before
any run exists. \emph{Adjoint forecast} is Theorem~\ref{thm:adjoint} run on the first seed of each
arm. \emph{Measured deviation} is given the realised across-seed deviation and is the ceiling.

\subsection{Result}

Table~\ref{tab:sizing} is the outcome. The adjoint forecast reaches
$\numSizeAdjointRuns$ runs against the ceiling's $\numSizeOracleRuns$: within resampling error of a
policy told the answer, at $\numSizeOracleGain\times$
$[\numSizeOracleLo, \numSizeOracleHi]$. Ordinary fixed-budget practice needs $\numSizeFixedRuns$
runs for the same reliability, $\numSizeFixedGain\times$ $[\numSizeFixedLo, \numSizeFixedHi]$ more.
Spending seeds to estimate the deviation instead of forecasting it costs
$\numSizeEmpiricalGain\times$ $[\numSizeEmpiricalLo, \numSizeEmpiricalHi]$, and most of that
penalty is structural: that policy cannot stop at one seed per arm whatever the lead looks like, so
its floor is four runs before any comparison is decided.

\begin{table}[t]
  \centering
  \small
  \begin{tabular}{lccc}
  \toprule
  policy & runs & vs.\ adjoint & $95\%$ \\
  \midrule
    fixed budget, no stopping rule & 30.41 & $9.57$ & $[7.24, 10.56]$ \\
    deviation from the seeds in hand & 6.52 & $2.13$ & $[1.98, 2.38]$ \\
    one deviation for every recipe & 6.04 & $1.93$ & $[1.45, 2.28]$ \\
    frozen forecast, before any run & 3.34 & $1.09$ & $[1.02, 1.14]$ \\
    adjoint forecast, from one run & 3.07 & $1.00$ & $[1.00, 1.00]$ \\
    measured deviation & 3.16 & $1.02$ & $[0.99, 1.05]$ \\
  \bottomrule
\end{tabular}
  \caption{Runs needed to hold the mistake rate at $\numSizeTarget\%$ over $\numSizePairs$
  unambiguous comparisons, both arms counted, with $95\%$ intervals from resampling the
  comparisons. The forecast is not better than knowing the deviation; it is not measurably worse,
  which is the point. Ordering follows forecast accuracy: the adjoint forecast is
  $\numSizeForecastError\times$ off at the median here, the frozen predictor
  is coarser, and a constant deviation discards the variation entirely.}
  \label{tab:sizing}
\end{table}

The ordering of the six tracks how accurately each predicts the deviation, which is the result
rather than a coincidence, and it is what makes the next question the operative one.

\subsection{How wrong the forecast may be}

On a pretrained model the forecast is accurate only in some regimes
(Section~\ref{sec:exp-support}). Whether that matters depends on how much error the decision
tolerates, so we corrupt the forecast on purpose and rerun the benchmark: a lognormal
multiplicative error of controlled median magnitude, and separately a systematic scale error.

The scale error does nothing. Over a range of $\numTolScaleRange$ orders of magnitude the runs
needed change by a factor of $\numTolScaleSpread$, because $z$ absorbs it exactly as
\eqref{eq:stop} says it should. A forecast does not need to be calibrated to be useful here, which
is worth stating plainly because it is the opposite of what is required for reporting an error bar.

Relative error is what costs. Runs needed rise from $\numTolClean$ at the measured accuracy to
$\numTolRunsOneHalf$ at a median error of $1.5\times$, $\numTolRunsDouble$ at $2\times$ and
$\numTolRunsTwoHalf$ at $2.5\times$, reaching fixed-budget's $\numTolFixed$ at about
$\numTolBreakeven\times$. So the tolerance is wide but not unlimited: a forecast wrong by less
than roughly a factor of three pays for itself, and one wrong by twenty-five, as
Section~\ref{sec:exp-real} reports for a corpus whose reward groups go unanimous, does not.

That break-even is the honest statement of scope. It is met comfortably on the exact models, it is
met by the accurate regime of the pretrained sweep, and it is not met by the others. The
contribution is a decision procedure together with the condition under which it beats what is
currently done, not a claim that the condition always holds.
